\subsubsection{恒等式的逐项展开与计算位宽}
单个配对项展开为
\[
(x_0+w_1)(x_1+w_0)=x_0x_1+x_0w_0+x_1w_1+w_0w_1.
\]
减去输入内部乘积 $x_0x_1$ 与权重内部乘积 $w_0w_1$ 后，恰好保留两个所需乘积。以 $x_0=3,x_1=-2,w_0=4,w_1=1$ 为算术示例，主项为 $4\times2=8$，行修正为 $-6$，列修正为 $4$，得到 $8-(-6)-4=10$，与直接计算 $3\times4+(-2)\times1=10$ 一致。该示例说明恒等关系，不代表某个测试样本或额外测试结果。

位宽需要按变换后的中间表达式确定。PE 端口保留 signed12；两输入相加扩展到 13 位，再产生 26 位乘积。32 个乘积累加使用 31 位，随后与扩展后的行因子、列因子及偏置一起进入 32 位修正通路。不能因为最终输出只有 12 位就提前截断主项，否则行列修正不再抵消，恒等式将失效。

\begin{table}[H]\centering\small
\caption{冻结数值通路的表示宽度与用途}
\begin{tabularx}{\linewidth}{p{34mm}p{24mm}Y}\toprule
量 & 表示宽度 & 位宽说明\\\midrule
输入 $X$、PE 端 $W$ & 12 位有符号 & W10 符号扩展至 12 位后送入阵列\\
交叉预加和 & 13 位有符号 & 保留两输入相加产生的最高有效位\\
单个配对乘积 & 26 位有符号 & 13 位乘 13 位的完整乘积\\
主项累加和 & 31 位有符号 & 覆盖 32 项累加所需的增长位\\
行、列修正项 & 29 位有符号 & 在共享存储及修正通路中扩展使用\\
偏置字段 & 22 位有符号 & 转换后数值放入原偏置字段\\
修正中间结果 & 32 位有符号 & 减去行、列因子并加入偏置\\
输出 & 12 位有符号 & 仅在最后执行舍入和饱和\\\bottomrule
\end{tabularx}
\end{table}

按端口允许的完整 signed12 范围作保守分析，预加和的绝对值不超过 4096，单项乘积绝对值不超过 $2^{24}$，32 项不超过 $2^{29}$；31 位有符号主项可容纳该保守界。实际 W10 的范围更小，但本版并未因此缩窄 PE 端口和权重物理条带。这里的位宽分析解释当前实现，不将尚未实施的位宽压缩计入面积收益。

\subsubsection{有效位与流水停顿}
乘法寄存器保存 product、first 和 last；累加寄存器保存同一输出位置的主项。若输出仍有效且下游不接收，累加级停止消费；乘法级在无法继续前推时保持数据及标记。first 随首个乘积进入累加级，以零作为基值，last 随末个乘积到达后发布结果。因此 first/last 是数值边界的一部分，不能独立于乘积推进。

\begin{figure}[H]\centering
\begin{tikzpicture}
\node[box,text width=32mm] (a) at (0,0) {预加与乘法\\输入握手};
\node[box,text width=35mm] (b) at (4.5,0) {product 寄存器\\first / last / valid};
\node[box,text width=35mm] (c) at (9.2,0) {累加寄存器\\末项发布 out\_valid};
\draw[arr](a)--(b);\draw[arr](b)--(c);
\draw[arr,accent](c.south)--++(0,-.8)-|node[pos=.3,below,note]{下游未接收时保持有效数据}(b.south);
\node[note,text width=105mm] at (4.5,-2.65) {复位清除控制有效位；新帧的 first 确定累加初值。\\数据寄存器的旧值不作为有效结果发布。};
\end{tikzpicture}
\caption{PE 两级存储与反压依赖示意。图中连接表示控制关系，不是采集波形。}
\end{figure}

该设计减少了宽数据寄存器统一复位的需求，但正确性依赖于首项标记、有效位和上层 32 对调度契约。当前整机回放与局部检查均保留这些约束，报告不把去除数据复位描述为可脱离有效位机制单独使用的优化。
